\section{Input dimension proportional to depth}
The same exponential-depth conclusion holds when input dimension grows linearly with depth. This is useful when comparing depth upper and lower bounds whose input dimensions are themselves proportional to depth.

\begin{theorem}[A uniform bound for $d=O(H)$]
\label{thm:linear-dimension-depth}
Fix $\kappa,C,B,\varepsilon,\delta$. There is a finite constant
$a=a(\kappa,C,B,\varepsilon,\delta)$ such that the conclusion of Theorem~\ref{thm:exp-depth-ReLU} holds for every $H\ge1$ and every
$d\le\kappa(H+1)$ whenever $n\ge e^{a(H+1)}$.
The normalized-energy certificate uses the sharper constants specified below. No other learner assumption changes.
\end{theorem}
\begin{proof}
Two coarse constants in the preceding proof can be kept exact. First,
\[
 \frac{D_N}{\omega L}\le\frac{D+H}{L}
       =1+\frac dL\le1+\kappa,
\]
so use $d_0=1+d/L$ in the tail and entropy bounds.
Second, the product calculation in the teacher-balancing step gives the per-layer bound
\[
 e=D^{2/L}C^{2H/L}B^2.
\]
This is bounded by a constant depending only on $\kappa,C,B$ when
$D=d+1\le(\kappa+1)L$: the quantity $L^{2/L}$ is uniformly bounded for $L\ge2$. If some layer energy vanishes, the represented function is zero and that case requires no balancing.

The large-width comparison ratio is still
\[
 \frac{D_{\rm nar}}{\omega L}
 \le\frac{C+(D+1)/L}{N^{\gamma-1}},
\]
whose numerator is bounded in terms of $C,\kappa$. Therefore
$T_N,K_n,\mathcal B$ remain $O_{\kappa,C,B}(1+\log n/L)$ when
$N^{\gamma-1}\ge L$.
The capacity coefficient satisfies $k_H^2=O_\kappa(L)$ because
$\log D=O_\kappa(\log L)$. Hence the same bound
$a_n^2\le c_{\kappa,\delta}2^L(1+t)^4L^5e^{-tL}$ holds.
Every step making the root excess small uniformly in $L$ is unchanged.

In the complementary branch $N<L$, both the input dimension and width are $O_\kappa(L)$. Parameter count is $O_\kappa(L^3)$ and pseudodimension is $O_\kappa(L^4\log(eL))$, as before. The exponential sample budget dominates this polynomial. The same certificate selection and confidence allocation complete the proof.
\end{proof}
